% =============================================================================
% sections/preliminary_test_plan.tex  -  Section 2: Preliminary Test Plan
%
% SCORING: 0--25 pts
%   Structure (0-5 pts):
%     - All requirements verified by test mentioned          (0-1 pts)
%     - Test ID, Name, Description provided                 (0-1 pts)
%     - Test setup described                                (0-1 pts)
%     - Pass/fail criteria and test status provided         (0-1 pts)
%     - Tests timeline provided                             (0-1 pts)
%   Technical excellence (0-20 pts)
% =============================================================================

\section{Preliminary Test Plan}
\label{sec:test_plan}

\subsection{Introduction}
\label{subsec:tp_intro}

This section presents the preliminary test plan for the \teamname\ rover and
drone system. The purpose of the test plan is to verify that all requirements
identified in Section~\ref{sec:moc} are satisfied before the competition.

\subsection{Test Overview}
\label{subsec:tp_overview}

Table~\ref{tab:test_summary} summarises all tests planned for the
verification campaign.

\begin{table}[H]
    \centering
    \caption{Test overview summary}
    \label{tab:test_summary}
    \begin{tabular}{lc}
        \toprule
        \textbf{Category}        & \textbf{Number of tests} \\
        \midrule
        Rover mechanical tests   & [X] \\
        Rover electrical tests   & [X] \\
        Drone mechanical tests   & [X] \\
        Drone electrical tests   & [X] \\
        Software / integration tests & [X] \\
        \midrule
        \textbf{Total}           & [X] \\
        \bottomrule
    \end{tabular}
\end{table}

\subsection{Test Descriptions}
\label{subsec:tp_descriptions}

Each test is described in Table~\ref{tab:test_plan}. The columns are:

\begin{description}[style=nextline]
    \item[Test ID] Unique identifier (e.g.\ T-01).
    \item[Test Name] Short, descriptive name.
    \item[Req. ID] Requirement(s) verified by this test (links to
          Section~\ref{sec:moc}).
    \item[Description] What the test involves and what is being measured.
    \item[Test Setup] Equipment, environment, and configuration needed.
    \item[Pass Criteria] Measurable condition(s) that must be met to pass.
    \item[Status] \textit{Planned} / \textit{In Progress} / \textit{Pass} /
          \textit{Fail}.
\end{description}

\begin{longtable}{C{0.8cm} L{2.2cm} C{1cm} L{2.8cm} L{2.8cm} L{2.2cm} C{1.2cm}}
    \caption{Preliminary Test Plan}
    \label{tab:test_plan} \\
    \toprule
    \textbf{ID} & \textbf{Test Name} & \textbf{Req.} & \textbf{Description}
        & \textbf{Setup} & \textbf{Pass Criteria} & \textbf{Status} \\
    \midrule
    \endfirsthead
    \multicolumn{7}{c}{\small\itshape (continued from previous page)} \\
    \toprule
    \textbf{ID} & \textbf{Test Name} & \textbf{Req.} & \textbf{Description}
        & \textbf{Setup} & \textbf{Pass Criteria} & \textbf{Status} \\
    \midrule
    \endhead
    \midrule
    \multicolumn{7}{r}{\small\itshape (continued on next page)} \\
    \endfoot
    \bottomrule
    \endlastfoot

    % ---- Rover tests --------------------------------------------------------
    T-01 & [Test name]  & R-01 & [Description of the test] & [Equipment / environment] & [Measurable criterion] & Planned \\
    T-02 & [Test name]  & R-02 & [Description of the test] & [Equipment / environment] & [Measurable criterion] & Planned \\
    % ---- Drone tests --------------------------------------------------------
    T-10 & [Test name]  & R-0X & [Description of the test] & [Equipment / environment] & [Measurable criterion] & Planned \\
    % Add all tests ...

\end{longtable}

\subsection{Test Timeline}
\label{subsec:tp_timeline}

Figure~\ref{fig:test_timeline} shows the planned schedule for the
verification campaign leading up to ERC~2026.

\begin{figure}[H]
    \centering
    % Replace with an actual Gantt chart or timeline figure
    \fbox{\parbox{0.9\textwidth}{\centering\vspace{2cm}
        [Insert test timeline / Gantt chart here]\vspace{2cm}}}
    \caption{Preliminary test campaign timeline}
    \label{fig:test_timeline}
\end{figure}

% Alternatively, use a table-based timeline:
\begin{table}[H]
    \centering
    \caption{Test timeline (quarters / months)}
    \label{tab:test_timeline}
    \begin{tabular}{L{3cm} C{1.5cm} C{1.5cm} C{1.5cm} C{1.5cm} C{1.5cm}}
        \toprule
        \textbf{Test}   & \textbf{Jan} & \textbf{Feb} & \textbf{Mar}
                        & \textbf{Apr} & \textbf{May} \\
        \midrule
        T-01 [name]     & $\bullet$    &              &  &  & \\
        T-02 [name]     &              & $\bullet$    &  &  & \\
        T-10 [name]     &              &              & $\bullet$ &  & \\
        % Add rows for each test
        \bottomrule
    \end{tabular}
\end{table}
